% Plantilla de actividad IMTC
\documentclass[spanish,letterpaper,oneside]{article}

\def\documenttitle {Inteligencia Artificial Redes Neuronales}
\def\documentsubtitle {Actividad 1 - inteligencia-artificial-redes-neuronales-imtc}
\def\documentsubject {Ingenieria en Mecatronica}
\def\documentauthor {Martin Jonathan de la Cruz}
\def\coursename {Inteligencia Artificial Redes Neuronales}
\def\coursecode {IMTC}
\def\universityname {Universidad Autonoma de Nuevo Leon}
\def\universityfaculty {Facultad de Ingenieria Mecanica y Electrica}
\def\universitydepartment {\coursename}
\def\universitydepartmentimage {UANL/assets-uanl/logo-uanl.png}
\def\universitydepartmentimagecfg {width=3.2cm}
\def\universitylocation {San Nicolas de los Garza, Nuevo Leon}
\def\authortable {
  \begin{tabular}{ll}
    \textbf{Alumno:} & Martin Jonathan de la Cruz \\
    Matricula: & UANL-IMTC \\
    Programa: & Ingenieria en Mecatronica \\
    Asignatura: & \coursename \\
    Codigo: & \coursecode \\
    Figura docente: & Nombre por definir \\
    & \\
    \multicolumn{2}{l}{Fecha de realizacion: \today} \\
    \multicolumn{2}{l}{\universitylocation}
  \end{tabular}
}

\input{template}

\begin{document}
\templatePortrait
\templatePagecfg
\onehalfspacing

\section{Actividad 1}
\textbf{Consigna tecnica:} Documentar y resolver el reto academico de la semana.

\subsection{Desarrollo}
Incluir evidencias, analisis y referencias.

\subsection{Cierre}
Registrar resultados, limitaciones y mejoras.

% Ciclo PDF: bibliographystyle omitido; la plantilla AulaTeX ya define el estilo bibliográfico
\bibliography{inteligencia-artificial-redes-neuronales-imtc}

% --- Ciclo A: refuerzo editorial materializado ---
\section*{Refuerzo editorial Ciclo A}
Este apartado consolida la mejora editorial incremental del nodo a partir del estándar maduro de AulaTeX. El refuerzo atiende brechas detectadas de bibliografía, citas, análisis propio, transferencia y cierre argumentado, sin sustituir la consigna original ni copiar contenido de los casos maestros.

\subsection*{Tesis de trabajo}
La actividad debe sostener una postura verificable: el producto solicitado no sólo organiza información, sino que permite interpretar el problema disciplinar, justificar decisiones y reconocer sus consecuencias académicas o profesionales.

\subsection*{Del registro descriptivo a la \textbf{interpretación disciplinar}}
Llamo \textbf{registro descriptivo} a la exposición que se limita a enumerar pasos ejecutados, capturas o resultados obtenidos, y \textbf{interpretación disciplinar} al ejercicio que explica qué significan esos resultados dentro del marco conceptual de la asignatura y qué decisiones los produjeron. El hallazgo central es que una entrega madura requiere transitar del primero a la segunda, pues el mismo resultado puede ser correcto en lo procedimental y, sin embargo, irrelevante en lo argumentativo si no se conecta con el problema que motiva la actividad \citep{programa-inteligencia-artificial-redes-neuronales-imtc,cicloARefuerzo02}. Por tanto, el desarrollo debe explicitar qué revela el producto construido, por qué resulta relevante para \coursename{} y con qué criterio se evalúa su pertinencia. Conviene subrayar, no obstante, que la interpretación no sustituye a la evidencia: la supone, la ordena y la somete a un juicio explícito \citep{cicloARefuerzo03}.

\subsection*{\textbf{Trazabilidad} documental de las decisiones}
Entiendo por \textbf{trazabilidad} la posibilidad de reconstruir, a partir del propio documento, el camino que va del planteamiento del reto a la conclusión: qué se asumió como dato, qué criterio guió cada elección y qué fuente respalda cada afirmación. En consecuencia, sostengo que la trazabilidad es la condición mínima de verificabilidad de la entrega, porque sin ella la postura propia se vuelve indistinguible de una opinión no fundamentada \citep{cicloARefuerzo04,cicloARefuerzo06}. Este criterio, además, articula el resto del refuerzo: la evidencia se vuelve auditable, la transferencia profesional se apoya en decisiones documentadas y el cierre puede recuperar la tesis sin introducir elementos ajenos al desarrollo.

\subsection*{\textbf{Evidencia}, criterio de evaluación y postura argumentada}
Sostengo que la evidencia solo adquiere valor académico cuando se articula con un criterio explícito de calidad y con una postura propia fundamentada; en consecuencia, no basta con mostrar resultados, sino que es necesario justificar decisiones y límites del análisis. Esta operación convierte la evidencia reunida en argumento académico verificable \citep{cicloARefuerzo03,cicloARefuerzo05}.

\subsection*{Transferencia profesional}
La transferencia consiste en vincular el producto con situaciones reales de desempeño: diagnóstico, toma de decisiones, comunicación de resultados, cumplimiento normativo, intervención educativa o resolución de problemas según el campo disciplinar. Así, la actividad adquiere valor formativo más allá de la plantilla.

\subsection*{Cierre argumentado}
En consecuencia, la calidad de la actividad depende de que el producto visible, las fuentes y la conclusión formen una secuencia coherente. La conclusión debe recuperar la tesis, indicar el principal aprendizaje y señalar una implicación práctica o conceptual.

% Brechas locales atendidas por Ciclo A: tesis, analisis_propio, transferencia, conclusion_argumentada, bibliografia_insuficiente, citas_insuficientes
% Reglas globales aplicadas:
% - Priorizar cierre de brecha recurrente: bibliografia_insuficiente (427 nodos).
% - Priorizar cierre de brecha recurrente: conclusion_argumentada (375 nodos).
% - Priorizar cierre de brecha recurrente: citas_insuficientes (328 nodos).
% - Priorizar cierre de brecha recurrente: analisis_propio (288 nodos).
% --- Fin Ciclo A ---


\subsection*{Citas de refuerzo Ciclo A}
La mejora editorial se vincula con la memoria documental disponible mediante las referencias \citep{programa-inteligencia-artificial-redes-neuronales-imtc,cicloARefuerzo02,cicloARefuerzo03,cicloARefuerzo04,cicloARefuerzo05,cicloARefuerzo06}. Estas citas deben revisarse en una etapa disciplinar fina para sustituir o complementar fuentes genéricas por literatura específica del tema.

\end{document}
